\chapter{The Synthesized Library: Architecture}
\label{ch:synthlib}

\chapterorient{Five parts have disassembled Palace into named pieces and shown
that every analysis is a composition of the same handful of them. This part puts
the pieces back---not as the C++ they came from, but as a small, clean
\emph{library} written in the calculus of Part~III: the simulator rendered as
code. This opening chapter lays out that library's \emph{architecture}. We meet
the three \emph{views} of the machine the book has been building (top-down,
bottom-up, and---new here---the implementation view), then the five modules the
library splits into, the dependency graph that orders them, the load-bearing rule
that keeps the code a \emph{view} of the spec rather than a second copy of it, and
the \texttt{\#extern} boundary where three opaque kernels stop. By the end you
will know the shape of the whole library; the next two chapters tour its modules
(\cref{ch:fivelibs}) and trace one driver through it as code
(\cref{ch:composing}).}

\section{Three views of one machine}
\label{sec:threeviews}

A complete simulator is a large object, and a large object is best understood from
more than one side. This book has, by now, walked around the machine and looked at
it from three. It is worth naming all three at once, because this part is the
third of them, and it only makes sense in relation to the other two.

The first view is \emph{top-down}: \emph{what does the machine do?} You stand at
the user-facing entry points---the six analyses an engineer actually invokes---and
read each as a coherent feature: the electrostatic capacitance extraction, the
driven frequency sweep, the eigenmode solve. Part~V took this view, devoting one
chapter to each modality and closing, in \cref{ch:fulltrace}, with a single driver
traced from its configuration file to its output plot. A driver, seen top-down, is
a \emph{composition root}: inputs are configuration, the output is the physical
product, and the body composes already-built vocabulary. The top-down view answers
``what is this analysis, end to end?''

The second view is \emph{bottom-up}: \emph{what is each piece?} You start at the
ground---cited Palace source---and climb, re-expressing each operation in a
cleaner vocabulary one rotation at a time, until you reach the small calculus of
combinators at the top (\cref{ch:operators}, \cref{ch:rotations}). Parts~II
through~IV took this view: the function spaces, the assembly fold, the Krylov
solvers, the multigrid preconditioners, the reduction verbs, each defined in
itself, with its signature, its laws, its place in the vocabulary. The bottom-up
view answers ``what is \code{fe\_assemble}, exactly, and where is it used?''

The third view is the one this part adds. It is the \emph{implementation view}:
\emph{what is the code?} Not what the machine does (top-down), not what each piece
is (bottom-up), but the actual library a programmer would read if the spec were
realized as software---the def bodies, written out, in the pseudo-language of
\cref{ch:language}, organized into modules, with every cross-reference resolving to
a real function in a real file. The implementation view answers ``show me the code
that realizes all of this.''

\begin{figure}[t]
\centering
\begin{tikzpicture}[font=\footnotesize, >=Stealth]
  % the central machine
  \node[stage, minimum width=26mm, minimum height=14mm, fill=simBgGray, draw=simGray,
        align=center] (m) at (0,0) {\textbf{the simulator}\\[1pt]\scriptsize one machine};
  % --- top-down view (from above) ---
  \node[product, align=center, text width=30mm] (td) at (0,2.7)
    {\textbf{top-down view}\\[1pt]\scriptsize \emph{what does it do?}\\6 driver features\\(Part~V, \cref{ch:fulltrace})};
  \draw[flow, simRust] (td.south) -- node[right, font=\scriptsize]{compose} (m.north);
  % --- bottom-up view (from below) ---
  \node[opbox, align=center, text width=30mm] (bu) at (0,-2.7)
    {\textbf{bottom-up view}\\[1pt]\scriptsize \emph{what is each piece?}\\vocabulary L0$\to$L4\\(Parts~II--IV)};
  \draw[flow, simTeal] (bu.north) -- node[right, font=\scriptsize]{decompose} (m.south);
  % --- implementation view (from the side: THIS PART) ---
  \node[stage, fill=simBgViolet, draw=simViolet, very thick, align=center, text width=32mm]
    (iv) at (5.6,0) {\textbf{implementation view}\\[1pt]\scriptsize \emph{what is the code?}\\the synthesized library\\(\textbf{this Part})};
  \draw[flow, simViolet] (iv.west) -- node[above, font=\scriptsize]{render} (m.east);
  % the three angles annotation
  \node[font=\scriptsize\itshape, text=simGray, align=center] at (0,-4.15)
    {three views of the \emph{same} machine---the simulator does not change,\\only the side we read it from};
\end{tikzpicture}
\caption[Three views of one machine]{The same simulator, read from three sides.
The \emph{top-down} view (above, rust) enters at the user-facing driver features
and asks what each analysis \emph{does}, reading it as a composition of vocabulary;
Part~V took this view. The \emph{bottom-up} view (below, teal) enters at cited
source and asks what each \emph{piece is}, climbing the vocabulary from L0 to L4;
Parts~II--IV took this view. The \emph{implementation view} (right, violet) is this
Part: it renders the same vocabulary as concrete library \emph{code}---the def
bodies a programmer would read. The machine is one object; the three views are
three angles on it, not three machines.}
\label{fig:threeviews}
\end{figure}

\Cref{fig:threeviews} sets the three side by side. They are not three machines;
they are three angles on one machine. The top-down view recomposes the vocabulary
into features; the bottom-up view decomposes the features into vocabulary; the
implementation view renders that same vocabulary as the code that realizes it. The
last is the natural endpoint of the journey: \cref{ch:bigidea} chose the
pure-function picture precisely because it is ``the form the accelerator wants''---a
library of pure tensor-in, tensor-out functions, ready to lower onto a GPU. This
part writes that library down.

\begin{keyidea}
The simulator can be read from three sides. \emph{Top-down} (Part~V): what does each
analysis \emph{do}---a driver as a composition root over the vocabulary.
\emph{Bottom-up} (Parts~II--IV): what is each \emph{piece}---the vocabulary defined
in itself, L0 to L4. \emph{The implementation view} (this Part): what is the
\emph{code}---the same vocabulary rendered as a small library of pure functions, the
bridge to a real GPU or accelerator backend. The synthesized library \emph{is} the
simulator written out as code; it is the destination \cref{ch:bigidea} pointed at
when it chose the pure-function picture for being ``the form the accelerator wants.''
\end{keyidea}

\subsection{Why a synthesized \emph{library}, and not a port}

A word on what we are and are not building. We are not porting Palace---not
translating its C++ line by line into another language. We are \emph{synthesizing}
a library: writing down, in one clean place, the code that the spec describes,
factored the way the spec factored it (combinator-primary, \cref{ch:combinators}),
with the optimization tricks of the original erased (\cref{ch:bigidea}) and the
pure-function reading made literal. The result is shorter than the original by a
large factor, because the original's bulk is mostly performance machinery and
boilerplate that the spec has already stripped away. What remains is the algorithm.

This is why the library is small enough to print. The composition tree of
\cref{ch:fulltrace}---the whole driven driver expanded to its matrix-free
leaves---was finite and shallow; rendered as code, each box in that tree is a few
lines, and the whole library is a few hundred lines across five modules. That
smallness is not a trick of presentation. It is the payoff of five parts of
disassembly: once you have found the handful of combinators that the whole
simulator is built from, writing them down as code is a short exercise.

\section{The five-library partition}
\label{sec:fivelib}

The library splits into five modules. The split is not arbitrary: three of the
five \emph{mirror} the three groups the calculus of \cref{ch:combinators} already
fell into, and the other two \emph{bracket} those three---one below, one above---to
make the whole stack come out in a clean order. \Cref{fig:bracket} is the hero
figure of the chapter: the five-library bracket, drawn as a stack.

\begin{figure}[p]
\centering
\resizebox{\textwidth}{!}{%
\begin{tikzpicture}[font=\footnotesize, >=Stealth]
  % ===== drivers cap (top) =====
  \node[product, minimum width=92mm, minimum height=11mm, align=center] (drv) at (0,5.2)
    {\textbf{\code{drivers}} \;\; \emph{the entry-point surfaces}\\[1pt]
     \scriptsize 6 sim drivers \;$\cdot$\; 6 output products \;$\cdot$\; the \code{lifecycle} ROOT \;\;(composes everything below)};
  % ===== three middle columns =====
  \node[stage, fill=simBgGreen, draw=simGreen, minimum width=28mm, minimum height=22mm,
        align=center] (iter) at (-3.1,2.6)
    {\textbf{\code{iteration}}\\[2pt]\scriptsize \code{iterate\_while}\\\code{krylov\_step}\\\code{chebyshev}\\[2pt]\itshape iteration \&\\step combinators};
  \node[stage, fill=simBgGold, draw=simGold, minimum width=28mm, minimum height=22mm,
        align=center] (data) at (0,2.6)
    {\textbf{\code{data-algebra}}\\[2pt]\scriptsize \code{linear\_combination}\\\code{inner\_product}\\\code{fe\_assemble} \dots\\[2pt]\itshape data-algebra\\combinators \& verbs};
  \node[stage, fill=simBgRust, draw=simRust, minimum width=28mm, minimum height=22mm,
        align=center] (coord) at (3.1,2.6)
    {\textbf{\code{coordination}}\\[2pt]\scriptsize \code{ksp\_solve}\\\code{solve\_family}\\\code{fold\_solve} \dots\\[2pt]\itshape outer-driver caps\\\& coordination};
  % ===== types slab (bottom) =====
  \node[stage, fill=simBgBlue, draw=simBlue, minimum width=92mm, minimum height=11mm,
        align=center] (types) at (0,-0.3)
    {\textbf{\code{types}} \;\; \emph{the shared cross-cutting records}\\[1pt]
     \scriptsize \code{IoData} \;$\cdot$\; \code{OpParams} \;$\cdot$\; \code{SimState} \;\;(every consumer rests on these)};
  % ===== brackets / arrows =====
  % drivers composes the three middle
  \draw[depedge] (drv.south west) ++(8mm,0) -- (iter.north);
  \draw[depedge] (drv.south)              -- (data.north);
  \draw[depedge] (drv.south east) ++(-8mm,0) -- (coord.north);
  % three middle rest on types
  \draw[depedge] (iter.south)  -- (types.north west) ++(8mm,0);
  \draw[depedge] (data.south)  -- (types.north);
  \draw[depedge] (coord.south) -- (types.north east) ++(-8mm,0);
  % the "mirrors the calculus" tag
  \node[font=\scriptsize\itshape, text=simGray, align=center] at (0,0.95)
    {the three middle libraries \emph{mirror} the three calculus doc-groups of \cref{ch:combinators}};
  % bracket labels on the side
  \node[font=\scriptsize\bfseries, text=simRust, anchor=west, align=left] at (5.2,5.2)
    {top bracket:\\\scriptsize composes down};
  \node[font=\scriptsize\bfseries, text=simBlue, anchor=west, align=left] at (5.2,-0.3)
    {bottom bracket:\\\scriptsize everything rests on it};
  % topological-order arrow
  \draw[->, simGray, thick] (-5.4,-0.3) -- (-5.4,5.2)
    node[midway, rotate=90, anchor=south, font=\scriptsize\itshape, text=simGray]
    {topological order: a def appears after everything it uses};
\end{tikzpicture}}
\caption[The five-library bracket]{The architecture of the synthesized library: a
five-module stack. The three middle libraries (\code{iteration},
\code{data-algebra}, \code{coordination}) \emph{mirror} the three doc-groups the
calculus of \cref{ch:combinators} already fell into---iteration and step
combinators, data-combining verbs, and outer-driver coordination. They are
\emph{bracketed} below by \code{types}, the genuinely shared record definitions
(\code{IoData}, \code{OpParams}, \code{SimState}) that every consumer rests on, and
above by \code{drivers}, the entry-point surfaces (the 6 drivers, 6 output
products, and the \code{lifecycle} ROOT) that \emph{compose} everything beneath
them. The bottom slab must come first and the top cap last because the stack is in
\emph{topological order}: a definition appears only after everything it uses, so
shared types precede their consumers and composing drivers come last. Read the
stack bottom to top and you are reading the library in the order you could actually
compile it.}
\label{fig:bracket}
\end{figure}

\subsection{The three middle libraries mirror the calculus}

{\sloppy
Recall \cref{ch:combinators}. The combinator vocabulary fell naturally into groups:
the \emph{iteration-structural} combinators that thread a carry through a loop or a
family (\code{iterate\_while} and its relatives), and the \emph{data-algebra}
combinators that combine tensors into tensors or scalars (\code{linear\_combination}
and \code{inner\_product}). The synthesized library keeps that grouping, refining it
into three modules:\par}

\begingroup\sloppy
\begin{description}[style=nextline, leftmargin=0pt, font=\normalfont\itshape]
  \item[\code{iteration} --- iteration \& step combinators.] The value-threaded loop
  primitive \code{iterate\_while} (\cref{ch:fold}) and the step kernels it folds:
  the Krylov step \code{krylov\_step}, the Chebyshev smoother \code{chebyshev}.
  These are the machines that drive one iterative solve to convergence.
  \item[\code{data-algebra} --- data-combining verbs.] The tensor-sum combinator
  \code{linear\_combination} and the scalar-producing \code{inner\_product}
  (\cref{ch:combinators}), together with the named verbs that the steps and
  reductions call: \code{fe\_assemble}, the reduction verbs \code{gram\_reduce} /
  \code{sparameter\_reduce} / the rest, the operator builders. These are the
  machines that combine data.
  \item[\code{coordination} --- outer-driver caps \& coordination.] The
  \code{Solve}-monadic outer driver \code{ksp\_solve} (\cref{ch:monad}) and the
  family combinators that drive it over a collection: \code{solve\_family},
  \code{frequency\_sweep}, \code{fold\_solve}, and the eigensolve cap
  \code{eigsolve}. These are the machines that coordinate many solves.
\end{description}
\endgroup

These three are the heart of the library, and they are exactly the three questions
\cref{ch:bigidea} taught us to ask of any computation: \emph{what operations}
(\code{data-algebra}), \emph{who owns the state} (the \code{Solve} monad, in
\code{coordination}), \emph{how is evolution coordinated} (the combinators, split
between \code{iteration} and \code{coordination}). The library's middle is the
calculus, modularized.

\subsection{The two brackets: \texttt{types} below, \texttt{drivers} above}

The three middle libraries do not stand alone. They are bracketed.

Beneath them sits \code{types}, the foundation. It holds only the records that are
\emph{genuinely shared}---used by two or more of the API groups above it. There are
three: \code{IoData}, the parsed configuration every solve reads (\cref{ch:records},
\cref{ch:strata}); \code{OpParams}, the read-only operator-internal parameters the
Krylov kernel and the coordination caps both close over; and \code{SimState}, the
solve-state record the iteration kernel evolves and the coordination caps thread.
A record used by only \emph{one} group does not live here---it is placed with that
group (the rule has a name, below). \code{types} is the cross-cutting remainder,
and it sits at the bottom because every consumer rests on it.

Above the three sits \code{drivers}, the cap. It holds the entry-point
surfaces---the same composition roots Part~V studied top-down: the six simulation
drivers, the six output products, and the \code{lifecycle} ROOT that dispatches
among them. Each is rendered as code that \emph{composes} the three middle
libraries. The electrostatic driver, written out, is a few lines calling
\code{fe\_assemble}, \code{solve\_family}, and \code{gram\_reduce}---exactly the
composition of \cref{ch:fulltrace}, now as a function. \code{drivers} sits at the
top because it uses everything below it and nothing uses it.

\begin{notationbox}
The three middle libraries take their names---\code{iteration},
\code{data-algebra}, \code{coordination}---directly from the three calculus
doc-groups of \cref{ch:combinators}; reading them is reading that chapter's two
combinator branches as modules. The two brackets, \code{types} and \code{drivers},
are the \emph{new} names this Part introduces, and they are bookkeeping names: the
shared-record floor and the entry-point cap. Five modules, three of them the
calculus you already know and two of them the brackets that order it.
\end{notationbox}

\subsection{Why a bracket, and why this order}

Why must \code{types} be first and \code{drivers} last? Because the library is
written in \emph{topological order}: a definition appears only after everything it
mentions has already been defined. This is the ordinary discipline of any
single-pass module---you cannot call a function before you have written it---and it
forces the bracket shape. A shared record like \code{SimState} is mentioned in the
signature of \code{krylov\_step} (which evolves it) and \code{ksp\_solve} (which
threads it); so \code{SimState} must be defined before both, which puts \code{types}
at the bottom. A driver like \code{electrostatic} mentions \code{fe\_assemble},
\code{solve\_family}, and \code{gram\_reduce}; so those must all be defined before
it, which puts \code{drivers} at the top. The brackets are not a stylistic choice;
they are where topological order \emph{forces} the shared floor and the composing
cap to go.

\begin{figure}[t]
\centering
\resizebox{\textwidth}{!}{%
\begin{tikzpicture}[font=\footnotesize, >=Stealth, node distance=4mm]
  % a linear "reading order" of defs, left to right, with dependency arcs pointing backward
  \def\dy{0}
  \node[stage, fill=simBgBlue, draw=simBlue] (t1) at (0,\dy) {\code{SimState}};
  \node[stage, fill=simBgBlue, draw=simBlue, right=of t1] (t2) {\code{OpParams}};
  \node[opbox, fill=simBgGreen, draw=simGreen, right=of t2] (iw) {\code{iterate\_while}};
  \node[opbox, fill=simBgGreen, draw=simGreen, right=of iw] (ks) {\code{krylov\_step}};
  \node[opbox, fill=simBgRust, draw=simRust, right=of ks] (kspv) {\code{ksp\_solve}};
  \node[opbox, fill=simBgRust, draw=simRust, right=of kspv] (sf) {\code{solve\_family}};
  \node[product, right=of sf] (el) {\code{electrostatic}};
  % the reading-order baseline
  \draw[->, simGray, very thick] (-1.0,-0.85) -- (12.6,-0.85)
    node[right, font=\scriptsize\itshape, text=simGray] {reading / compile order};
  % drop ticks
  \foreach \n in {t1,t2,iw,ks,kspv,sf,el}{
    \draw[simGray!50] (\n.south) -- (\n.south |- 0,-0.85);}
  % dependency arcs (each points BACKWARD to something already defined): u -> v means u uses v
  \draw[depedge] (ks.north) to[out=120,in=60] (iw.north);          % krylov_step uses iterate_while
  \draw[depedge] (ks.north) to[out=140,in=40] (t1.north);          % ...and SimState
  \draw[depedge] (kspv.north) to[out=120,in=60] (ks.north);        % ksp_solve folds krylov_step
  \draw[depedge] (kspv.north) to[out=150,in=30] (t2.north);        % ...and OpParams
  \draw[depedge] (sf.north) to[out=120,in=60] (kspv.north);        % solve_family maps ksp_solve
  \draw[depedge] (el.north) to[out=120,in=60] (sf.north);          % electrostatic uses solve_family
  % annotation: every arc points LEFT (backward)
  \node[font=\scriptsize\itshape, text=simBlue, align=center] at (6.0,2.0)
    {every dependency arc points \emph{backward}: a def appears only after everything it uses};
\end{tikzpicture}}
\caption[Topological definition order]{Topological order, read along one line. The
library's definitions are laid out left to right in the order a single-pass reader
(or compiler) meets them; the lower arc is that reading order. Each curved arc above
is a \code{depends-on} edge---\code{krylov\_step} uses \code{iterate\_while} and
\code{SimState}; \code{ksp\_solve} folds \code{krylov\_step} and reads
\code{OpParams}; \code{solve\_family} maps \code{ksp\_solve}; the
\code{electrostatic} driver composes \code{solve\_family}. The defining property is
that \emph{every arc points backward}: a definition appears only after everything it
uses has already appeared. This is exactly why the shared \code{types}
(\code{SimState}, \code{OpParams}) come first and the composing driver comes last---a
forward-pointing arc would be a use-before-definition, which topological order
forbids. The five-library bracket of \cref{fig:bracket} is this one line, folded into
a stack.}
\label{fig:topo}
\end{figure}

\Cref{fig:topo} draws the order as a single line of defs with their dependency arcs,
and makes the defining property visible at a glance: every arc points
\emph{backward}. A forward-pointing arc would name something not yet defined---a
use-before-definition---and topological order is precisely the absence of such arcs.
The bracket of \cref{fig:bracket} is this line folded into a stack: the floor is its
left end, the cap its right.

\begin{keyidea}
The library is five modules: a \code{types} floor, three middle libraries
(\code{iteration}, \code{data-algebra}, \code{coordination}) that \emph{mirror} the
three calculus doc-groups of \cref{ch:combinators}, and a \code{drivers} cap that
\emph{composes} them into the entry-point surfaces of Part~V. The shape is a
bracket because the library is in \emph{topological order}---a def appears after
everything it uses---so the genuinely shared records must come first (the floor)
and the composing drivers must come last (the cap). Read bottom to top, the stack
is the order you could compile it in.
\end{keyidea}

\subsection{The type-placement rule}
\label{sec:typeplacement}

One subtlety of the floor deserves its own statement, because it is the rule that
keeps \code{types} small. \code{types} holds only records shared across \emph{two or
more} API groups. A record used by exactly \emph{one} group is not cross-cutting,
and putting it in the floor would force the reader to hold an unfamiliar type in
mind long before reaching the code that uses it. So a single-group record is placed
\emph{immediately before its own group}, bundled with its \emph{utility} API---its
constructors, field accessors, and trivial projections, the small namespace
intrinsic to the type. Its \emph{consumer} methods---the substantive operators that
read the record to do real work---stay in the group proper, after the type.

The Krylov workspace \code{Krylov} is used only by the iteration kernels, so it
lives at the head of \code{iteration}, not in the floor. The eigensolve state
\code{EigState} is used only by \code{eigsolve}, so it lives at the head of
\code{coordination}. The per-driver configuration records are thin projections of
\code{IoData} used only by their own driver, so they live at the head of
\code{drivers}. Only the truly cross-cutting three---\code{IoData}, \code{OpParams},
\code{SimState}---earn a place in the shared floor. Topological order still governs
throughout: a type and its utility API always precede their consumers, whether the
type sits in the floor or at the head of its group.

\section{The dependency graph and its two edge classes}
\label{sec:depgraph}

The bracket of \cref{fig:bracket} is a picture of \emph{dependencies}: who uses
whom. To make the picture precise, we draw the dependency graph explicitly and
distinguish its two kinds of edge, because the distinction is what makes the
ordering well-defined.

A \emph{dependency} between two defs comes in two flavors. A \textbf{\code{depends-on}}
edge is a real build dependency: def $u$ \code{depends-on} def $v$ when $u$'s body
\emph{calls} $v$, so $v$ must be defined before $u$ can be. These edges are
\emph{blocking}---they constrain the order, and they are exactly the edges
topological order must respect. \code{ksp\_solve} \code{depends-on}
\code{krylov\_step} (it folds it); \code{electrostatic} \code{depends-on}
\code{fe\_assemble}, \code{solve\_family}, and \code{gram\_reduce} (it composes
them); every consumer \code{depends-on} the \code{types} it reads.

A \textbf{\code{reference}} edge is navigational only: $u$ \code{reference}s $v$
when $u$ \emph{points the reader at} $v$ without calling it---a cross-link, a ``see
also,'' a ``this is the same idea as.'' These edges are \emph{free}: they do not
constrain the order and they impose no build dependency. When the \code{driven}
driver's documentation notes that its port modes ``come from the boundary-mode
analysis,'' that is a \code{reference}, not a \code{depends-on}: \code{driven} does
not \emph{call} \code{boundary\_mode}; it consumes a product that happens to be
produced elsewhere, and the cross-link is a courtesy to the reader, not a wire in
the build graph.

\begin{figure}[t]
\centering
\resizebox{\textwidth}{!}{%
\begin{tikzpicture}[font=\footnotesize, >=Stealth, node distance=10mm and 16mm]
  % the stack, as nodes
  \node[product, minimum width=24mm] (drv) at (0,3.6) {\code{drivers}};
  \node[stage, fill=simBgGreen, draw=simGreen, minimum width=20mm] (it) at (-3,1.8) {\code{iteration}};
  \node[stage, fill=simBgGold,  draw=simGold,  minimum width=22mm] (da) at (0,1.8) {\code{data-algebra}};
  \node[stage, fill=simBgRust,  draw=simRust,  minimum width=22mm] (co) at (3,1.8) {\code{coordination}};
  \node[stage, fill=simBgBlue,  draw=simBlue,  minimum width=24mm] (ty) at (0,0) {\code{types}};
  % depends-on edges (solid blue): drivers -> three middle -> types
  \draw[depedge] (drv) -- (it);
  \draw[depedge] (drv) -- (da);
  \draw[depedge] (drv) -- (co);
  \draw[depedge] (it) -- (ty);
  \draw[depedge] (da) -- (ty);
  \draw[depedge] (co) -- (ty);
  % an internal depends-on: coordination -> iteration (ksp_solve folds krylov_step)
  \draw[depedge] (co) -- node[above, font=\scriptsize, pos=0.45]{folds} (it);
  % coordination -> data-algebra (drives the verbs)
  % a reference edge (dotted gray): a driver cross-links a sibling driver's product
  \draw[refedge] (drv.east) .. controls +(2.4,-0.4) and +(1.6,1.4) ..
    node[right, font=\scriptsize, pos=0.5, text=simGray]{\emph{reference} (sibling product)} (co.east);
  % legend (sample edges drawn directly, with adjacent labels)
  \draw[depedge] (-6.6,3.6) -- (-6.0,3.6);
  \node[anchor=west, font=\scriptsize] at (-5.9,3.6) {\code{depends-on}};
  \node[anchor=west, font=\scriptsize, text=simGray] at (-6.6,3.15) {blocking; constrains order};
  \draw[refedge] (-6.6,2.4) -- (-6.0,2.4);
  \node[anchor=west, font=\scriptsize] at (-5.9,2.4) {\code{reference}};
  \node[anchor=west, font=\scriptsize, text=simGray] at (-6.6,1.95) {navigational; free};
\end{tikzpicture}}
\caption[The dependency graph and its two edge classes]{The library's dependency
graph, with its two edge classes distinguished. A \code{depends-on} edge (solid
blue) is a real build dependency---$u$ \code{depends-on} $v$ when $u$'s body calls
$v$---so it \emph{constrains the order}: the arrows all point downward, from
\code{drivers} through the three middle libraries to \code{types}, and an internal
arrow records that \code{coordination} folds \code{iteration}'s kernels. These are
the edges topological order respects. A \code{reference} edge (dotted gray) is
merely navigational---a cross-link to a sibling's product, here a driver pointing
at another driver's output it consumes---and is \emph{free}: it imposes no build
dependency and constrains nothing. The stack \code{types} $\leftarrow$
\{\code{iteration}, \code{data-algebra}, \code{coordination}\} $\leftarrow$
\code{drivers} is exactly the \code{depends-on} order.}
\label{fig:depgraph}
\end{figure}

\Cref{fig:depgraph} draws the graph. The \code{depends-on} edges (solid) all run
downward: \code{drivers} \code{depends-on} the three middle libraries, each of which
\code{depends-on} \code{types}, with the one internal cross-edge that
\code{coordination} folds \code{iteration}'s step kernels. Follow only the solid
edges and you recover the bracket stack of \cref{fig:bracket}; that is the whole
content of ``topological order''---the solid arrows never point up. The
\code{reference} edges (dotted) are the free cross-links the reader navigates by;
they are allowed to point any direction, because they bind nothing.

\begin{pitfall}
The distinction between \code{depends-on} and \code{reference} is not pedantry; it
is what keeps the dependency graph \emph{well-founded}. If we treated every
cross-link as a build dependency, the graph would tangle: the driven driver
``references'' the boundary-mode driver, the boundary-mode driver shares vocabulary
with the eigenmode driver, and soon every module appears to depend on every other
and no topological order exists. By reserving \code{depends-on} for actual calls and
demoting cross-links to free \code{reference} edges, the build graph stays a clean
downward stack while the reader still gets every helpful pointer. A cross-link that
is \emph{not} a call must be a \code{reference}, or the order breaks.
\end{pitfall}

\section{The implementation view versus the semantics}
\label{sec:viewvssem}

We come now to the load-bearing distinction of the whole part---the one rule that
makes the synthesized library a \emph{view} of the specification rather than a
second, rivalrous copy of it. State it plainly: \emph{the library renders the def
bodies; it does not restate the semantics.}

Every operator in this book has a definitional home---a chapter where its
signature, its algebraic laws, its meaning, and (for a record) its field schema are
stated \emph{once} and authoritatively. \code{linear\_combination}'s laws live in
\cref{ch:combinators}; \code{ksp\_solve}'s convergence semantics live in the Krylov
chapters; \code{SimState}'s field schema lives with the records of
\cref{ch:records}. The synthesized library does \emph{not} repeat any of that. It
renders the \emph{code}---the def body, the concrete expression that computes the
result---and for everything else it \emph{links} to the authoritative home. A
rendered \code{linear\_combination} in the library is a few lines of foldl; it does
not re-derive the concatenation-homomorphism law, it points at \cref{ch:combinators}
where that law is proved.

\begin{figure}[t]
\centering
\begin{tikzpicture}[font=\footnotesize, >=Stealth]
  % LEFT: the semantics home
  \node[stage, fill=simBgGold, draw=simGold, minimum width=34mm, minimum height=22mm,
        align=center] (sem) at (-3.4,0)
    {\textbf{the semantics home}\\[2pt]\scriptsize (the operator chapter,\\\cref{ch:combinators};\\the calculus, Part~III)\\[3pt]
     \footnotesize signature\\laws\\meaning\\record schema};
  \node[font=\scriptsize\itshape, text=simGold, anchor=north] at (-3.4,-1.7)
    {stated ONCE, authoritative};
  % RIGHT: the rendered code
  \node[stage, fill=simBgViolet, draw=simViolet, minimum width=34mm, minimum height=22mm,
        align=center] (code) at (3.4,0)
    {\textbf{the synthesized def}\\[2pt]\scriptsize (the library chapter,\\this Part)\\[4pt]
     \footnotesize the def \emph{body}:\\the concrete code\\that computes it};
  \node[font=\scriptsize\itshape, text=simViolet, anchor=north] at (3.4,-1.7)
    {the implementation VIEW};
  % the link (reference) from code to semantics
  \draw[refedge, very thick] (code.west) -- node[above, font=\scriptsize]{\emph{links to}}
    node[below, font=\scriptsize, text=simGray]{(does not restate)} (sem.east);
  % the "no back-copy" forbidden arrow
  \draw[-{Stealth[length=2mm]}, simRust, dashed, thick] (sem.north east) ++(0,0.2)
    .. controls +(1.5,1.0) and +(-1.5,1.0) .. (code.north west) ++(0,0.2);
  \node[font=\scriptsize\itshape, text=simRust] at (0,2.05) {$\times$ no second copy of the laws};
\end{tikzpicture}
\caption[Implementation view versus semantics]{The load-bearing distinction. The
\emph{semantics} of an operator---its signature, its algebraic laws, its meaning,
and (for a record) its field schema---live \emph{once}, authoritatively, in its
home chapter (left, gold): \code{linear\_combination}'s laws in
\cref{ch:combinators}, a record's schema with \cref{ch:records}. The synthesized
library (right, violet) renders only the def \emph{body}---the concrete code that
computes the result---and \emph{links} to the home for everything else (solid
arrow). It does \emph{not} copy the laws across (the forbidden dashed arrow): the
library is a \emph{view} of the spec, like generated code is a view of its source,
not a second authoritative statement of it. This is what keeps the code and the spec
from drifting into two rivalrous copies.}
\label{fig:viewvssem}
\end{figure}

\Cref{fig:viewvssem} draws the rule. Why insist on it? Because a specification and
its implementation, if both state the laws, will eventually \emph{disagree}---someone
edits one and forgets the other, and now there are two answers to ``what does this
operator mean.'' By making the library a pure \emph{view}---bodies plus links, no
restated laws---we guarantee there is exactly one authoritative statement of every
semantic fact, and the code is checked \emph{against} it rather than competing with
it. This is the same relationship generated code has to its source: the generated
form is real, runnable, and useful to read, but it is not where the meaning
\emph{lives}; the meaning lives in the source it was generated from. The synthesized
library is generated code in exactly this sense, and the calculus of Part~III is its
source.

\begin{keyidea}
The synthesized library is a \emph{view}, not a second copy. It renders the def
\emph{bodies}---the concrete code---and \emph{links} to each operator's home chapter
for the signature, laws, meaning, and record schema, which are stated \emph{once and
only once} there. It never restates the semantics. This is the discipline that keeps
the implementation and the specification from drifting into two rivalrous truths:
there is one authoritative statement of every fact, and the code is a checkable view
of it, exactly as generated code is a view of its source. The implementation view is
parallel to the top-down feature spine---another angle on the one machine, not a new
machine.
\end{keyidea}

\section{The \texttt{\#extern} boundary: where the library stops}
\label{sec:extern}

There is one more thing the library must be honest about: its \emph{edges}. Not
every operation the simulator performs is ours to write. Three of them bottom out in
an opaque external kernel---a specialized library Palace calls but does not
itself implement---and the synthesized code must mark those boundaries rather than
fabricate a body for them. The mark is \code{\#extern}.

When a def's implementation lives behind an opaque library boundary, the library
renders its \emph{type signature} and then, in place of a body, writes
\code{\#extern NAME}. This says, exactly: \emph{here is the contract---the input and
output types---and here is where our code stops and the external kernel takes over;
we do not fabricate what we do not own.} Three kernels in the whole simulator earn
an \code{\#extern}:

\begin{description}[style=nextline, leftmargin=0pt, font=\normalfont\itshape]
  \item[The libCEED element-quadrature kernel.] Inside \code{fe\_assemble}
  (\cref{ch:assembly}), the element-local quadrature contraction is performed by
  libCEED, the matrix-free kernel library (\cref{ch:matrixfree}). The assembly fold
  is ours; the per-element kernel it calls is \code{\#extern}.
  \item[The SLEPc eigensolve loop.] Inside \code{eigsolve} (\cref{ch:krylovschur}),
  the eigenvalue iteration itself runs inside SLEPc---Palace writes no loop, it
  hands the whole problem across and waits. The cap is ours; the iteration is
  \code{\#extern}.
  \item[The ODE time-step.] Inside \code{fold\_solve} (\cref{ch:time}), the transient
  analysis's per-step time advance bottoms out in an MFEM time-integrator (an
  implicit Runge--Kutta or generalized-$\alpha$ step). The fold that threads the
  state is ours; the step it folds is \code{\#extern}.
\end{description}

\begin{figure}[t]
\centering
\begin{tikzpicture}[font=\footnotesize, >=Stealth]
  % the def: type signature (ours) + extern body (opaque) + the impl link
  \node[stage, minimum width=40mm, minimum height=8mm, align=center] (sig) at (0,2.4)
    {\code{eigen\_iterate :: OpParams -> ...}\\[1pt]\scriptsize the type signature (ours: the contract)};
  \node[extern, minimum width=40mm, minimum height=9mm, align=center] (ext) at (0,1.0)
    {\code{\#extern eigen\_iterate}\\[1pt]\scriptsize the opaque kernel (SLEPc EPS loop);\\our code stops here};
  \draw[flow] (sig) -- (ext);
  % the kernel-API label
  \node[font=\scriptsize\bfseries, text=simViolet, anchor=west] at (3.2,1.0)
    {= the kernel-\emph{API}};
  \node[font=\scriptsize\itshape, text=simGray, anchor=west, align=left] at (3.2,0.55)
    {the boundary / contract\\(what the spine calls)};
  % the separate kernel-IMPL node, linked by realizes-kernel-api (dashed ref)
  \node[stage, fill=simBgGreen, draw=simGreen, minimum width=46mm, minimum height=12mm,
        align=center] (impl) at (0,-1.4)
    {\textbf{the kernel-\emph{implementation}}\\[1pt]\scriptsize a SEPARATE constructive realization\\in OUR vocabulary: Lanczos / Arnoldi /\\Krylov--Schur (\cref{ch:krylovschur})};
  \draw[refedge, very thick] (ext.south) --
    node[right, font=\scriptsize, pos=0.5]{\code{realizes-kernel-api}}
    node[left, font=\scriptsize, text=simGray, pos=0.5, align=right]{(reviewable;\\not a call)} (impl.north);
\end{tikzpicture}
\caption[The \texttt{\#extern} boundary and the kernel-API/impl split]{The
\code{\#extern} boundary, and the distinction it does \emph{not} erase. An opaque
kernel renders as its \emph{type signature} (top: the contract, ours to state)
followed by \code{\#extern NAME} (middle, violet dashed) in place of a body: our
code stops here and the external library---here SLEPc's eigenvalue loop---takes
over. This \code{\#extern} surface is the \emph{kernel-API}: the boundary and
contract the simulator calls across. It is \emph{not} the last word. Separately, the
book gives a \emph{kernel-implementation} (bottom, green): a constructive realization
of the same kernel in our own vocabulary---for the eigensolve, the Lanczos / Arnoldi
/ Krylov--Schur iteration of \cref{ch:krylovschur}. The two are linked by a
\code{realizes-kernel-api} edge (dashed): a \emph{reviewable correspondence}, not a
build call---a reader checks that the from-our-primitives version matches the opaque
contract. The \code{\#extern} marks where \emph{the library} stops; the kernel-impl
is where \emph{the book} keeps going.}
\label{fig:extern}
\end{figure}

\Cref{fig:extern} draws the boundary and the crucial thing it does \emph{not} do:
it does not abandon the question. An \code{\#extern} surface is the
\emph{kernel-API}---the contract the simulator calls across, the opaque library
boundary. But for each of these kernels the book \emph{separately} provides a
\emph{kernel-implementation}: a constructive realization of the same operation in
our own vocabulary, built from primitives we already own. The SLEPc eigensolve loop
is \code{\#extern} in the library, but \cref{ch:krylovschur} reconstructs what that
loop does---the Lanczos, Arnoldi, and Krylov--Schur iterations---in our calculus, so
the curious reader can read both the black-box contract and a glass-box realization
of it. The two are tied together by a \code{realizes-kernel-api} link: a
\emph{reviewable correspondence}, not a build call. The implementation does not
\code{depends-on} the API and the API does not call the implementation; rather, a
reader checks that the from-our-primitives version computes what the opaque contract
promises.

\begin{notationbox}
\code{\#extern NAME} is not a stub and not a failure---it is an honest edge. It
appears \emph{after} the type signature, never before: the library always states the
\emph{contract} (the input and output types) even where it does not own the
\emph{body}. Read \code{\#extern} as ``the contract is here; the implementation is
external; and---for these three---the book gives a constructive realization
separately, linked by \code{realizes-kernel-api}.'' Three kernels carry it: libCEED
quadrature, the SLEPc eigensolve loop, the ODE time-step.
\end{notationbox}

\begin{keyidea}
The library marks its edges with \code{\#extern}. Where an operation bottoms out in
an opaque external kernel---libCEED quadrature inside \code{fe\_assemble}, the SLEPc
eigensolve loop inside \code{eigsolve}, the ODE step inside \code{fold\_solve}---the
library renders the type signature and then writes \code{\#extern NAME} in place of
a body: it states the contract but does not fabricate what it does not own. The
\code{\#extern} surface is the \emph{kernel-API}. Distinct from it is the
\emph{kernel-implementation}: a separate, constructive realization of the same kernel
in our own vocabulary (e.g. the Krylov--Schur iteration of \cref{ch:krylovschur} for
the eigensolve), tied to the API by a \emph{reviewable} \code{realizes-kernel-api}
link, not a build call. The \code{\#extern} is where the \emph{library} stops; the
kernel-impl is where the \emph{book} keeps going.
\end{keyidea}

\section{Worked examples}

\begin{workedexample}{Reading the topological order, and placing \texttt{ksp\_solve}}{topo}
We make ``topological order'' concrete by deriving why \code{types} must come first
and \code{lifecycle} last, and then placing one specific operator,
\code{ksp\_solve}, in its module with its dependencies all pointing downward.

\emph{Why \code{types} first.} Pick any record in the floor, say \code{SimState}.
It is mentioned in the \emph{signature} of \code{krylov\_step} (whose monadic effect
\emph{is} the \code{SimState} transition) and in the signature of \code{ksp\_solve}
(which returns a \code{SimState}). Topological order forbids mentioning a name before
it is defined, so \code{SimState} must precede both \code{krylov\_step} (in
\code{iteration}) and \code{ksp\_solve} (in \code{coordination}). The same holds for
\code{OpParams} (read by the kernel and the caps) and \code{IoData} (read by every
driver). Three records, each used by two or more groups, each forced to precede its
users---so the cross-cutting records form a floor, and the floor is \code{types}.

\emph{Why \code{lifecycle} last.} The \code{lifecycle} ROOT \emph{dispatches} on the
problem type to one of the six drivers and calls it; so it mentions all six driver
names, which must already be defined. The drivers in turn mention the middle-library
verbs, which mention the floor. Nothing mentions \code{lifecycle}---it is the entry
point, the thing the user invokes, so no other def calls it. A def that uses
everything and is used by nothing is, by definition, topologically \emph{last}.

\emph{Placing \code{ksp\_solve}.} The preconditioned Krylov solve cap is an
outer-driver coordination machine: it runs the outer loop, classifies termination,
and projects the terminal state. So it belongs in \code{coordination}. Read its
dependencies:
\begin{lstlisting}[style=l4style]
ksp_solve :: OpParams -> Inputs -> SimState
ksp_solve op inp = execState (solve_loop op inp) (initial_state inp)
  where
    solve_loop op inp = do
      o <- one_cycle op inp
      unless (done o) (solve_loop op inp)
    one_cycle op inp = do
      s <- get
      let (kn, _) = iterate_while (krylov_step op (fresh_krylov op inp s)) cont
      modify (\s -> s { x = s.x `plus` (kn.basis `applyBasis` kn.y) })
      pure (classify kn op s)
\end{lstlisting}
Every name on the right-hand side is in a \emph{lower or equal} stratum:
\code{OpParams}, \code{SimState}, \code{Inputs} are in \code{types} (the floor);
\code{krylov\_step} and \code{iterate\_while} are in \code{iteration} (a middle
sibling, defined before \code{coordination} because \code{coordination} folds it);
\code{execState}, \code{get}, \code{modify} are the \code{Solve}-monad utilities
defined at the head of \code{coordination} itself, before \code{ksp\_solve}. No
dependency points \emph{up}: \code{ksp\_solve} never mentions a driver, never
mentions \code{lifecycle}. Its \code{depends-on} edges all run downward---\code{types}
and \code{iteration}---so it sits cleanly in \code{coordination}, above its
dependencies and below the \code{drivers} that will fold it into
\code{solve\_family}. That downward-only property is exactly what topological order
\emph{is}, checked at one node.
\end{workedexample}

\begin{workedexample}{The \texttt{\#extern} boundary inside \texttt{fe\_assemble}}{externassemble}
We render the assembly fold with its opaque kernel marked, and read off precisely
what the boundary means---contrasting it with the kernel-\emph{implementation} that
\cref{ch:matrixfree} provides.

The assembly verb \code{fe\_assemble} (\cref{ch:assembly}) folds a list of
weak-form terms over the mesh, accumulating each term's element-local contributions
into the operator. The \emph{fold} is ours---it is the ordinary
\code{linear\_combination}-shaped accumulation of \cref{ch:combinators}. But the
innermost step, evaluating one term's quadrature contraction on one element, is
performed by libCEED, the matrix-free kernel library. So the synthesized library
renders the fold and marks the kernel:
\begin{lstlisting}[style=l4style]
-- the assembly fold: ours. the per-element quadrature kernel: external.
fe_assemble :: FESpace -> [WeakTerm] -> LinOp
fe_assemble space terms =
  foldl (\op term -> op `plus_op` assemble_term space term) zero_op terms
  where
    -- one weak-form term -> its operator contribution, via the libCEED kernel.
    -- the CONTRACT (types) is ours; the BODY is the opaque matrix-free kernel.
    assemble_term :: FESpace -> WeakTerm -> LinOp
    #extern assemble_term
\end{lstlisting}
Read the two parts. The body of \code{fe\_assemble} is a plain \code{foldl}: it is
our code, the assembly fold of \cref{ch:assembly}, accumulating one operator
contribution per term. But \code{assemble\_term}---the per-element quadrature
contraction that turns one weak-form term into its matrix-free operator
action---renders as \code{\#extern} after its signature. That is the boundary: we
\emph{call into} libCEED's element-quadrature kernel, and the body is not ours to
fabricate. We state the contract (a \code{FESpace} and a \code{WeakTerm} in, a
\code{LinOp} out) and stop.

\emph{The contrast with the kernel-impl.} The \code{\#extern} is not the end of the
story for this kernel either. \Cref{ch:matrixfree} gives the
kernel-\emph{implementation}: a constructive realization of exactly this element
quadrature as a sequence of tensor contractions---gather the element DoFs, apply the
basis-evaluation tensor, contract against the quadrature weights and the geometry
factors, apply the basis-transpose, scatter back. That realization is built from
tensor primitives we already own, and it \emph{realizes-kernel-api} the
\code{\#extern assemble\_term} contract above: a reviewable correspondence, checked
by reading both, not a build call. So the same quadrature kernel appears twice in the
book, deliberately: once as the opaque contract the assembly fold calls
(\code{\#extern}, here), and once as the from-our-primitives tensor-contraction
implementation (\cref{ch:matrixfree}). The library renders the call site; the book
supplies the glass box beside it. That double rendering---API where the library
stops, impl where the book continues---is exactly the kernel-API/kernel-impl
distinction of \cref{sec:extern}, made concrete on the assembly kernel.
\end{workedexample}

\section{What the next two chapters do}
\label{sec:whatnext}

We have the architecture. \Cref{fig:bracket} is the floor plan: five modules, a
\code{types} floor, three middle libraries mirroring the calculus, a \code{drivers}
cap, all in topological order, with the \code{depends-on} edges running downward and
the \code{reference} cross-links free. We know the library is a \emph{view}---bodies
and links, never a second copy of the laws---and we know where it stops, at the three
\code{\#extern} kernels whose contracts it states but whose bodies it does not own.

The next two chapters fill the plan in. \Cref{ch:fivelibs} tours the five libraries
one at a time, opening each module and reading its rendered defs---the actual code,
module by module, from the \code{types} floor up through the \code{drivers} cap.
\Cref{ch:composing} then takes one driver---the electrostatic capacitance extraction,
the cleanest of the six---and composes it end to end \emph{in code}, the rendered
form of the very composition tree \cref{ch:fulltrace} expanded as a diagram. Together
they turn this architecture into the running library it describes, and
\cref{ch:wholemachine} closes the book by standing back from the whole machine and
asking where it goes next---onto the GPU and accelerator backends the pure-function
picture was chosen for in the first place.

\begin{keyidea}
\textbf{The whole simulator is one small library of composed functions---the
implementation view, mirroring the calculus.} Five modules: a \code{types} floor of
shared records, three middle libraries (\code{iteration}, \code{data-algebra},
\code{coordination}) that \emph{are} the three calculus doc-groups of
\cref{ch:combinators} rendered as code, and a \code{drivers} cap that composes them
into the entry-point surfaces of Part~V---all in topological order. It renders the
def \emph{bodies} and \emph{links} to the spec for the laws (a view, not a copy), and
marks its three opaque edges with \code{\#extern}. A hundred-thousand-line C++
simulator becomes a few hundred lines across five small modules, because once you
have found the handful of combinators it is built from, writing them as code is
short. That smallness is the whole book's thesis, finally rendered as the program it
always was.
\end{keyidea}

\summarysection
This chapter opens Part~VI, which renders the whole specification as a small
\emph{synthesized library}---the simulator written out as code in the calculus of
Part~III. We first named the three \emph{views} of the one machine: the
\emph{top-down} view of Part~V (what each analysis \emph{does}, a driver as a
composition root), the \emph{bottom-up} view of Parts~II--IV (what each \emph{piece
is}, the vocabulary L0$\to$L4), and---new here---the \emph{implementation view} (what
the \emph{code} is), the destination \cref{ch:bigidea} pointed at when it chose the
pure-function picture for matching the accelerator backend. The library splits into
\emph{five modules}: a \code{types} floor (the genuinely shared records
\code{IoData}, \code{OpParams}, \code{SimState}); three middle libraries
(\code{iteration}, \code{data-algebra}, \code{coordination}) that \emph{mirror} the
three calculus doc-groups of \cref{ch:combinators}; and a \code{drivers} cap that
\emph{composes} them into Part~V's entry-point surfaces. The bracket shape is forced
by \emph{topological order}---a def appears after everything it uses---so shared types
come first and composing drivers last; the type-placement rule keeps the floor small
by placing single-group records with their group. The dependency graph has two edge
classes: \code{depends-on} (a real call; blocking; constrains the order; the downward
arrows of the stack) and \code{reference} (navigational; free; the cross-links),
whose distinction keeps the graph well-founded. The load-bearing discipline is that
the library is a \emph{view}: it renders the def \emph{bodies} and \emph{links} to
each operator's home for the laws, meaning, and record schema---stated once
there---never restating them, exactly as generated code is a view of its source.
Finally, three opaque kernels (libCEED quadrature in \code{fe\_assemble}, the SLEPc
eigensolve loop in \code{eigsolve}, the ODE step in \code{fold\_solve}) render as
\code{\#extern NAME} after their type signature: the library states the contract but
does not fabricate the body. Each \code{\#extern} kernel-\emph{API} has a separate
kernel-\emph{implementation}---a constructive realization in our own vocabulary
(e.g. \cref{ch:krylovschur} for the eigensolve)---tied to it by a \emph{reviewable}
\code{realizes-kernel-api} link. The payoff: a hundred-thousand-line simulator becomes
a few hundred lines across five small modules; \cref{ch:fivelibs} tours them and
\cref{ch:composing} composes one driver as code.

\furtherreading
The pseudo-language the library is rendered in---Haskell-style type signatures over
pure functions, with do-notation for the state monad---is the notation of functional
programming; the standard reference for the language that crystallized it is the
\emph{Haskell} report and Peyton Jones's account of its design~\citep{peytonjones2003},
and the modular, type-directed style the five-library split follows is its everyday
practice. The \code{\#extern} libCEED quadrature kernel is documented in the CEED
project's matrix-free finite-element interface~\citep{ceed2021}, whose element-local
evaluation is exactly the contract \code{assemble\_term} renders as external and
\cref{ch:matrixfree} reconstructs as tensor contractions. The simulator whose surface
this library synthesizes is Palace~\citep{palace2023}; the synthesized form here is a
clean-room \emph{view} of that codebase's algorithm, not a transcription of its C++,
and the correspondence between the two---rendered body against cited source---is the
reviewable link the whole book has been building. The three-view framing and the
combinator-primary organization the middle libraries mirror were developed in
\cref{ch:bigidea} and \cref{ch:combinators}; this chapter only assembles them into an
architecture.

\exercisessection
\begin{enumerate}[nosep]
  \item \textbf{Name the three views.} State, in one sentence each, the question
  that the top-down view, the bottom-up view, and the implementation view answer
  about the simulator. Which Parts of the book take each view? Why is the
  implementation view the natural \emph{last} one, given the choice
  \cref{ch:bigidea} made about how to describe a computation?
  \item \textbf{Place a record.} The Krylov workspace \code{Krylov} is read only by
  the iteration kernels; the configuration record \code{IoData} is read by every
  driver and the coordination caps. Using the type-placement rule
  (\S\ref{sec:typeplacement}), say which module each belongs in and where within
  it, and explain why \code{IoData} earns a place in the \code{types} floor while
  \code{Krylov} does not.
  \item \textbf{Two edge classes.} The driven driver consumes port modes produced by
  the boundary-mode analysis but does not call \code{boundary\_mode}. Is the link
  from \code{driven} to \code{boundary\_mode} a \code{depends-on} edge or a
  \code{reference} edge? Justify your answer from the definitions, and explain what
  would go wrong with the topological order if every such cross-link were treated as
  a \code{depends-on}.
  \item \textbf{The view discipline.} The library's rendered \code{linear\_combination}
  is a few lines of \code{foldl} and contains \emph{no} statement of the
  concatenation-homomorphism law. Where does that law live, and why is it
  \emph{deliberately} absent from the rendered code? State the failure mode the
  ``view, not a copy'' discipline (\S\ref{sec:viewvssem}) is designed to prevent.
  \item \textbf{Mark the externs.} Of the three \code{\#extern} kernels (libCEED
  quadrature, the SLEPc eigensolve loop, the ODE time-step), say which module's def
  each appears inside, and what part of that def is \emph{ours} versus external. For
  one of them, explain in a sentence what the corresponding kernel-\emph{implementation}
  provides and which chapter gives it.
  \item \textbf{Topological check at a node.} Take \code{solve\_family}, which is
  \lstinline[style=l4style]|map (\inp -> ksp_solve op inp) rhss|. List every name on
  its right-hand side and assign each to a module; confirm that none of them sits in
  a module \emph{above} \code{coordination} (where \code{solve\_family} lives). What
  single fact about \code{solve\_family}'s dependencies makes it legal to place it in
  \code{coordination}?
  \item \textbf{Why the library is short.} The original Palace is roughly
  $10^5$ lines of C++; the synthesized library is a few hundred lines across five
  modules. Drawing on \cref{ch:bigidea}'s distinction between transparent
  optimization tricks and the base algebra, give two reasons the synthesized form is
  so much smaller, and name one thing that is \emph{deliberately not} reproduced in
  it. Is the library a \emph{port} of Palace? Explain the difference.
\end{enumerate}
